
\section{Preliminaries}\label{sec:prelims}

$\N$ denotes the set of non-negative integers and $[n]\coloneqq\set{1,\ldots,n}$ for $n\in\N$.
For a set $V$, we write $\binom V2$ for the set of unordered pairs $uv$ of distinct elements of $V$, and for $X,Y\subseteq V$ we put $XY\coloneqq\setof{xy}{x\in X,\ y\in Y,\ x\ne y}$. A graph $G$ consists of a finite set $V$ of \emph{vertices} and a set $E\subseteq {V\choose 2}$ of \emph{edges}, denoted by $V(G)$ and $E(G)$.
The (open) \emph{neighborhood} of a vertex $v\in V(G)$ is the set $N^G(v)\coloneqq\setof{w\in V(G)}{vw\in E(G)}$.
% lax begin Lax185574.GraphBalls
For $r\in \N$, denote $\Ball^G_r(v)\coloneqq\setof{w\in V(G)}{\dist^G(v,w)\le r}$, where $\dist^G(\cdot,\cdot)$ is the shortest-path distance in~$G$.
These notations extend to sets $A,B\subseteq V(G)$ with $\Ball_r^G(A) \coloneqq \bigcup_{a\in A} \Ball_r^G(a)$ and $\dist^G(A,B)\coloneqq\min_{a\in A,\,b\in B}\dist^G(a,b)$, and to tuples of vertices as expected.
% lax end

Let $G$ be a graph with vertex set $V$ and edge set $E$. By $\overline G$ we denote the edge-complement of $G$, which has  vertex set $V$ and edge set ${V\choose 2}-E$.
The subgraph of $G$ \emph{induced} by a set $X\subseteq V$ is the graph $G[X]$ with vertex set $X$ and edge set $E\cap {X\choose 2}$. For two disjoint sets $X,Y\subseteq V$,
by $G[X,Y]$ we denote the bipartite graph with parts $X$ and $Y$ and edge set $XY\cap E$.

A pair of sets $X,Y\subseteq V(G)$ is \emph{complete} if $XY\subseteq E(G)$ and \emph{anticomplete} if $XY\subseteq E(\overline G)$. The pair is \emph{homogeneous} if it is complete or anticomplete.

A \emph{graph class} is a set $\CC$ of graphs.

% lax begin Lax185574.Partitions
\paragraph{Partitions.}
For two partitions $\PP,\QQ$ of the same set we write $\PP\preceq \QQ$ to denote that $\PP$ refines $\QQ$, i.e., each part of $\PP$ is contained in some part of $\QQ$.
If $\PP$ is a partition of $V$ and $X\subseteq V$, denote
$X/\PP\coloneqq\{P\in\PP:P\cap X\ne\emptyset\}$.
The part of $\PP$ containing an element $v$ is denoted by $\PP(v)$.
Further, $\equiv_\PP$ is the equivalence relation on $V$ whose classes are the parts of $\PP$.
% lax end

\paragraph{Asymptotic notation.}
For a possibly empty list of parameters $k$,
 by $O_k(\cdot)$ we denote
 an anonymous function $f\from \R_{\ge 0}\to \R_{\ge 0}$ (depending only on the parameters~$k$) which is upper bounded by $x\mapsto a\cdot x+b$, for some $a,b\in\N$.

Similarly, we write $\poly_k(\cdot)$ for an anonymous function $f\from \R_{\ge 0}\to \R_{\ge 0}$ which is upper bounded by some polynomial%
;
in other words, $g(n)\le \poly_k(n)$ has the same implications as $g(n)\le (n+2)^{O_k(1)}$.
We use $\widehat O_k(n^c)$ for a bound $O_{k,\eps}(n^{c+\eps})$ for every $\eps>0$; thus it hides a subpolynomial factor that may depend on the parameters~$k$.

% lax begin Lax185574.MonadicDependence
\paragraph{Monadic dependence.} We view graphs as relational structures in the signature with one (symmetric, irreflexive) binary relation $E$.
For $k\in\N$, a \emph{$k$-colored graph} is a graph equipped with $k$ unary predicates $U_1,\ldots,U_k$.
A \emph{$k$-coloring} $G^+$ of a graph $G$ assigns interpretations of these predicates on $V(G)$; the predicates need not be disjoint.
A \emph{coloring} of $G$ is a $k$-coloring for some $k\in\N$.

A graph class $\CC$ is \emph{monadically independent} if there are $k\in\N$ and a first-order formula $\varphi(x,y)$ in the signature of $k$-colored graphs such that  every finite graph $H$ can be obtained from some graph $G\in\CC$ as follows: first choose a $k$-coloring $G^+$ of $G$, then consider the graph $\varphi(G^+)$ with vertex set $V(G)$ and edges $uv\in\binom{V(G)}2$ such that $G^+\models\varphi(u,v)\vee\varphi(v,u)$; finally take some induced subgraph of $\varphi(G^+)$, obtaining a graph isomorphic to $H$.
A graph class $\CC$ is \emph{monadically dependent} if it is not monadically independent.
% lax end

% lax begin Lax185574.VCDimension
\paragraph{Neighborhood complexity.}
The \emph{neighborhood complexity} of a graph $G$ is the function $\pi_G\from\N\to\N$ defined by:
\[
    \pi_G(k) \coloneqq
    \max_{\substack{S\subseteq V(G)\\|S|\le k}}
    \left|\{N^G(v)\cap S:v\in V(G)\}\right|,
\]
 which counts the maximum number of distinct neighborhoods of vertices in $G$ restricted to a set of size at most~$k$.

The \emph{VC-dimension} of a graph $G$ is the maximal size of a set $A\subseteq V(G)$ such that $\setof{N(v)\cap A}{v\in V(G)}=2^A$. The Sauer--Shelah lemma~\cite{Sauer72,She72} implies the following.
% lax end
% lax begin Lax185574.NeighborhoodComplexityBound
% lax begin Lax185574Proofs.NeighborhoodComplexityBound.neighborhoodComplexity_le
\begin{lemma}\label{lem:ss}
  If $G$ has VC-dimension $d$ then
  $$\pi_G(k)\le (k+1)^d\quad\text{for all }k\in\N.$$
\end{lemma}
% lax end
% lax end
% lax begin Lax185574Proofs.BoundedVCClasses.exists_vcDim_le
The following is easy to see from the definition of monadic dependence; see~\cite[Observation~3.2]{md-abmw} for a proof.
% lax end
% lax begin Lax185574.BoundedVCClasses
\begin{observation}\label{obs:md-vc}
For every monadically dependent graph class $\CC$ there is some $d\in\N$
such that every graph $G\in\CC$ has VC-dimension at most $d$.
\end{observation}
% lax end

% lax begin Lax185574.FirstOrderTypes
\paragraph{Types and parameters.}
For a structure $A$ and a set $W$ of its elements,
let $A_W$ denote the structure obtained by expanding $A$  with constants representing each element of $W$.
A formula $\phi(x_1,\ldots,x_k)$ \emph{with parameters} from $W$ is a formula in the signature of the structure $A_W$.
For $\bar v\in V(A)^k$ we write $A\models\phi(\bar v)$ to denote $A_W\models\phi(\bar v)$.

For $k,q\in\N$
and a tuple $\bar v\in V(A)^k$,
define
\emph{the type of $\bar v$ of quantifier rank $q$ over $W$}, denoted
$$\tp^A_q(\bar v/W),$$ as the set of all formulas $\phi(x_1,\ldots,x_k)$ with parameters from $W$ and of quantifier rank at most $q$ such that
$A\models \phi(\bar v)$.

For two sets $X,Y\subseteq V(A)$,
denote $$\Types^A_{k,q}(X/Y)\coloneqq \setof{\tp^A_q(\bar v/Y)}{\bar v\in X^k}.$$
We omit $/Y$ from the notation if $Y=\emptyset$.
% lax begin Lax185574.AtomicTypes
The \emph{atomic type} of $\bar v$ over $Y$ is defined as $\atp^A(\bar v/Y) \coloneqq \tp_0^A(\bar v/Y)$ and we set $\ATypes^A(X/Y)\coloneqq \Types^A_{1,0}(X/Y)$.

Note that for a graph $G$ and vertex set $S\subseteq V(G)$, two vertices $u,v\in V(G)$ have the same atomic type over $S$ if and only if either $u=v$, or $u,v\notin S$
and $N(u)\cap S=N(v)\cap S$.
The \emph{partition into atomic types}, denoted $\PP_S$, is the corresponding partition of
$V(G)$, which partitions $S$ into singletons, and vertices in $V(G)-S$ according to their neighborhood in $S$.
% lax end

% lax end
% lax begin Lax185574.GaifmanGraph
\paragraph{Gaifman graph, distances, balls.}
For a structure $A$ with domain $V$ equipped with unary and binary relations, the \emph{Gaifman graph} of $A$ is the graph
with vertex set $V$ and edge set consisting of all pairs $uv$ such that
$(u,v)$ or $(v,u)$ is in some binary relation of $A$.
We write $\dist^A(u,v)$ to denote the distance between two vertices in the Gaifman graph of $A$, and $\Ball^A_r(v)$  for the set of vertices at distance at most $r$ from $v$ in the Gaifman graph of $A$.

% lax end
% lax begin Lax185574.Flips
\paragraph{Flips and flip metrics.}
For a graph $G$ and partition $\PP$ of $V\coloneqq V(G)$, a \emph{$\PP$-flip} of $G$ is a graph $H$ with $V(H)=V$ such that for every pair $X,Y\in \PP$
(possibly $X=Y$), either
$H[X,Y]=G[X,Y]$ or $H[X,Y]=\overline G[X,Y]$.
If $|\PP| \leq s$, then we call $H$ an \emph{$s$-flip of $G$}.

Define the associated \emph{flip metric} and its balls by:
\begin{align*}
  \dist_\PP(u,v)&\coloneqq
  \max\{\dist^H(u,v):H\text{ is a $\PP$-flip of }G\},
  \\
  \Ball^\PP_r(v)&\coloneqq\{u:\dist_\PP(u,v)\le r\}=\bigcap\setof{\Ball^H_r(v)}{\text{$H$ is a $\PP$-flip of $G$}}.
\end{align*}

As a shorthand, for a graph $G$ and a set $S\subseteq V(G)$, a $\PP_S$-flip of $G$ is called an \emph{$S$-flip}, and we set $\dist_S^G \coloneqq \dist_{\PP_S}$ and $\Ball_r^S \coloneqq \Ball_r^{\PP_S}$.
% lax end

% lax begin Lax185574.ConstructionWidth
\paragraph{Partial graphs and merge-width.}
Partial graphs, construction sequences, their radius-$r$ width, and the radius-$r$ merge-width of a graph are defined in \Cref{sec:tech-overview} (see also~\cite{mw-stoc}).
Moreover, for a partial graph $H$ with partition $\PP$ and a set $W\subseteq V(H)$, the induced partial graph $H[W]$ is the partial graph with vertex set $W$,
edges $E(H)\cap {W\choose 2}$, non-edges
$N(H)\cap {W\choose 2}$, and partition $\setof{P\cap W}{P\in \PP}-\set{\emptyset}$.
% lax end
